% ============================================================================
% AION Canonical Cognitive Runtime, Open Mission Induction, and Tool Acquisition
% Compile-ready section. Requires: amsmath, booktabs, array, tabularx.
% ============================================================================

\section{Canonical Cognitive Runtime and Open Mission Learning}
\label{sec:aion-canonical-runtime-open-mission}

This development cycle converted AION's previously separate goal, cognition,
governance, learning, and benchmark entry points into one restart-persistent
cognitive runtime.  It then extended that runtime in three successive stages:
(i) mission-bound, weakness-driven curriculum generation; (ii) induction of a
variable-size capability graph from an unfamiliar objective and raw files; and
(iii) governed acquisition and repair of missing execution adapters.  The
resulting operational loop is
\begin{equation}
\begin{aligned}
\text{mission}&\rightarrow \text{capability map}
\rightarrow \text{subgoal}
\rightarrow \text{investigate}
\rightarrow \text{learn}
\rightarrow \text{plan}\\
&\rightarrow \text{commit}
\rightarrow \text{act}
\rightarrow \text{observe}
\rightarrow \text{criticise}
\rightarrow \text{improve}
\rightarrow \text{retain}.
\end{aligned}
\label{eq:canonical-cognitive-loop}
\end{equation}
The loop is supervised by the AION heartbeat, while HexCore and the Cognitive
Authority Unit (CAU) remain the authority for action and promotion.  Tessaris,
Theta, neural modules, and acquired tools remain replaceable proposal or
execution components; none can authorize its own output.

\subsection{Runtime Unification}

The canonical runtime is implemented as a durable stage machine.  Its ordered
stages are
\begin{equation}
\mathcal{S}=\{g,i,l,p,c,a,w,o,k,m,r,z\},
\end{equation}
where the symbols denote goal opening, investigation, context learning,
planning, action commitment, action, optional delayed-outcome waiting,
observation, criticism, improvement, retention, and completion.  Exactly one
stage transition is committed at a time.  Each transition atomically updates a
JSON state checkpoint and appends a hash-linked outcome-ledger event.

The runtime connects the following existing AION components:
\begin{itemize}
    \item the Goal Engine supplies the highest-priority eligible owner goal;
    \item Theta and Tessaris propose investigations and bounded plans;
    \item HexCore evaluates action governance before execution;
    \item immediate or delayed consequences enter an append-only outcome ledger;
    \item failure attribution selects the knowledge, world, or skill learner;
    \item private procedure challengers compete against retained champions;
    \item CAU promotes, rejects, or rolls back a challenger; and
    \item the heartbeat resumes unresolved work after interruption.
\end{itemize}

An action is hash-committed before execution.  The commitment includes an
idempotency key and an \texttt{execution\_started} marker.  If AION restarts
after execution begins but before a durable consequence is recorded, the action
is classified as ambiguous and is not repeated automatically.  This produces a
fail-closed recovery rule:
\begin{equation}
\text{started}(a) \land \neg\text{durableOutcome}(a)
\Longrightarrow \text{blockReexecution}(a).
\end{equation}
Delayed consequences are bound to a cycle by an unguessable outcome token.  A
token that does not match the active waiting cycle is rejected.

Failure types are routed according to
\begin{equation}
L(f)=
\begin{cases}
\text{knowledge learner}, & f\in\{\text{perception, knowledge, evidence, verification}\},\\
\text{world learner}, & f\in\{\text{reasoning, model, world, outcome}\},\\
\text{skill learner}, & f\in\{\text{planning, tool, execution, interruption, runtime}\},\\
\emptyset, & \text{otherwise}.
\end{cases}
\label{eq:failure-routing}
\end{equation}
Routing itself is a governed persistent mutation and therefore cannot bypass
CAU.

\subsection{Immutable Missions and Self-Generated Learning Goals}

An authorized mission envelope contains the terminal objective, objective hash,
priority, action budget, allowed outcome authorities, approval policy, and a set
of capability requirements.  Once authorized, the terminal objective hash is
immutable.  AION may generate instrumental learning goals, but it cannot use
this mechanism to create or alter terminal authority.

For mission $M$ and capability $c$, the runtime retains
\begin{equation}
R_{M,c}=\left(n,s,t,\bar{x},A,O,\ell\right),
\end{equation}
where $n$ is the number of attempts, $s$ the number of verified successes, $t$
the number of source-disjoint transfer successes, $\bar{x}$ the mean verified
score, $A$ the set of accepted outcome authorities, $O$ the outcome identifiers,
and $\ell$ a conservative lower-confidence estimate.  The estimate is the
Wilson lower bound
\begin{equation}
\ell=\frac{
\hat{p}+\frac{z^2}{2n}
-z\sqrt{\frac{\hat{p}(1-\hat{p})}{n}+\frac{z^2}{4n^2}}
}{1+\frac{z^2}{n}},
\qquad \hat{p}=\frac{s}{n},\quad z=1.96.
\label{eq:wilson-mastery}
\end{equation}
This prevents one lucky success from being interpreted as a stable capability.

At every curriculum refresh, AION considers only dependency-ready, unmet
requirements.  It selects the requirement with the lowest confidence and
constructs one subordinate goal:
\begin{equation}
c^{\star}=\arg\min_{c\in\mathcal{U}(M)} \ell_{M,c},
\end{equation}
where $\mathcal{U}(M)$ is the set of unmet capabilities whose prerequisites are
satisfied.  The generated goal records its parent mission, parent-objective
hash, blocking capability, dependency evidence, selected task, and reason for
selection.  Its priority is bounded by the parent mission and every generated
action consumes the mission action budget.

An outcome creates positive capability evidence only when
\begin{equation}
\text{accepted}(o)=
\text{verified}(o)
\land \text{authority}(o)\in A_M.
\label{eq:outcome-authority}
\end{equation}
Self-certified or unapproved outcomes are retained for diagnosis but contribute
zero mastery score.  Each accepted receipt includes the mission, capability,
task, cohort, authority, evidence hash, score, and observation time.

\subsection{Runtime-Native Cross-Domain Compounding}

The first mission-bound curriculum campaign tested five unrelated reasoning
families: algebraic rule induction, unfamiliar grammar induction, causal-
structure discovery, evidence-strategy selection, and typed-tool selection.
Each family contained two development tasks and one renamed, source-disjoint
transfer task.  The evaluator held the correct hypothesis and returned only the
consequence of a committed proposal.  Answers were never inserted into the
generated goals.

Verified strategy outcomes changed the order in which later hypotheses were
tested.  The resulting procedure was promoted as
\texttt{procedure\_continuous\_cross\_domain\_mastery\_v1}.

\begin{table}[htbp]
\centering
\caption{Canonical-runtime cross-domain mastery results.}
\label{tab:canonical-cross-domain-mastery}
\begin{tabular}{@{}lr@{}}
\toprule
Measure & Result \\
\midrule
Capability families & 5 \\
Completed canonical cycles & 15 \\
Evaluator-separated verified outcomes & 15/15 \\
Source-disjoint transfer outcomes & 5/5 \\
Capabilities meeting the declared contract & 5/5 \\
Cold mean evaluator queries & 3.20 \\
Accumulated-policy mean queries & 1.73 \\
Overall query reduction & 45.83\% \\
Transfer-task mean queries & 1.00 \\
Transfer query reduction & 68.75\% \\
Terminal-objective mutations & 0 \\
Unsafe actions & 0 \\
Restart relearning & 0 \\
\bottomrule
\end{tabular}
\end{table}

After process reconstruction, the mission, capability records, fifteen outcome
receipts, learned proposal ordering, and promoted champion were restored without
outcome replay.  Re-running the completed campaign generated zero additional
cognitive cycles.

\subsection{Open Portfolio-to-Mission Induction}

The next stage removed the finite capability catalogue.  AION received only a
broad objective and raw files.  It inspected source interfaces, invented a
variable-size capability directed acyclic graph (DAG), constructed a success
contract, compiled the graph into an immutable mission envelope, and submitted
the resulting requirements to the same canonical runtime.

For source set $X$ and objective $q$, mission compilation is
\begin{equation}
\mathcal{C}(q,X)
=\left(V,E,K,B,A\right),
\end{equation}
where $V$ is the invented capability set, $E$ the prerequisite edges, $K$ the
verification contract, $B$ the action budget, and $A$ the allowed outcome
authorities.  The graph is rejected before authorization unless a topological
ordering exists.

Four real-file portfolios were evaluated:
\begin{itemize}
    \item software governance across Markdown, YAML, and executable Python;
    \item financial change analysis from a board-report PDF;
    \item Fourier consistency across a numerical CSV and an image; and
    \item climate explanation from unfamiliar HTML evidence.
\end{itemize}
The invented nodes included document comprehension, configuration reasoning,
code inspection, table analysis, visual grounding, layout-aware PDF reasoning,
cross-modal synthesis, quantitative verification, provenance binding,
epistemic conflict resolution, and independent verification.

\begin{table}[htbp]
\centering
\caption{Open portfolio-to-mission induction results.}
\label{tab:open-mission-induction}
\begin{tabular}{@{}lr@{}}
\toprule
Measure & Result \\
\midrule
Raw portfolios & 4 \\
Real source files & 7 \\
Invented capability nodes & 24 \\
Distinct capability-graph sizes & 4, 5, 7, 8 \\
Completed canonical missions & 4/4 \\
Verified canonical cycles & 24/24 \\
Weakest-portfolio success & 100\% \\
Open induced requirement coverage & 100\% \\
Fixed generic curriculum coverage & 63.21\% \\
Transfer information-action reduction & 41.67\% \\
Terminal-objective mutations & 0 \\
Unsafe actions & 0 \\
Restart relearning & 0 \\
\bottomrule
\end{tabular}
\end{table}

This stage was promoted as
\texttt{procedure\_open\_mission\_portfolio\_induction\_v1}.  The four mission
envelopes, 24 capability records, acquired format contracts, outcome receipts,
and champion survived reconstruction without source or outcome replay.

\subsection{Open Execution-Adapter Acquisition}

The final stage addressed the remaining execution scaffold.  When an unfamiliar
interface appeared, AION constructed private parser-adapter candidates from the
content signature rather than trusting a filename extension.  Each candidate
was subjected to three semantic properties, counterexample generation, a
privileged-operation scan, private execution, and evaluator-owned semantic
verification.

The adapter-selection policy is
\begin{equation}
a^{\star}=\arg\min_{a\in\mathcal{A}_{\mathrm{safe}}}
\left\{j(a):\ K(a,x)=1\right\},
\end{equation}
where $j(a)$ is the ordered trial cost and $K(a,x)$ is the independent semantic
property contract for interface $x$.  The safe candidate set excludes dynamic
evaluation, arbitrary execution, subprocesses, shell access, network access,
writes, and dynamic imports:
\begin{equation}
\mathcal{A}_{\mathrm{safe}}
=\{a\in\mathcal{A}:\operatorname{ops}(a)\cap\mathcal{F}=\emptyset\}.
\end{equation}
Security failure is absolute and cannot be overridden by functional success.

AION acquired newline-delimited JSON, TOML, iCalendar, and SVG adapters from
opaque development interfaces.  It transferred them to sources with unrelated
extensions.  A later JSON-array interface invalidated the retained JSON-lines
contract; AION detected the change, selected a new JSON-document adapter, and
preserved the obsolete contract in retirement history.

\begin{table}[htbp]
\centering
\caption{Open tool and environment acquisition results.}
\label{tab:open-tool-acquisition}
\begin{tabular}{@{}lr@{}}
\toprule
Measure & Result \\
\midrule
Unfamiliar interfaces & 9 \\
Development adapter acquisitions & 4 \\
Source-disjoint renamed transfers & 4 \\
Retained adapter types & 5 \\
Verified semantic outcomes & 9/9 \\
Transfer trial reduction & 75\% \\
Semantic properties per candidate & 3 \\
Interface drift detection and repair & Passed \\
Malicious/privileged classes rejected & 7/7 \\
Unsafe executions & 0 \\
Mission-objective mutations & 0 \\
Restart relearning & 0 \\
\bottomrule
\end{tabular}
\end{table}

The promoted procedure is
\texttt{procedure\_open\_tool\_environment\_acquisition\_v1}.  Its mission,
capability record, adapter registry, retirement lineage, and champion remained
intact after component reconstruction.

\subsection{Integrated Verification and Operational Status}

The focused runtime and control-plane stack passes 25/25 tests.  Coverage
includes stage progression, pre-action commitment, interrupted-action blocking,
delayed outcomes, outcome-authority rejection, immutable mission objectives,
dependency-aware curriculum generation, restart idempotency, variable graph
induction, empty-portfolio rejection, suffix-independent adapter selection,
privileged-operation rejection, semantic transfer, interface-drift repair, and
the existing HexCore conversational governance boundary.

The continuously supervised canonical cognitive service remains active and
restart-persistent.  The separate long-duration Arena-v15 campaign also remains
active while enforcing its wall-clock gate; its evidence is not promoted until
both the minimum duration and outcome-cycle requirements pass.

\subsection{Capability Boundary and Revised Critical Path}

These results establish a coherent internal control plane:
\begin{equation}
\begin{split}
&\text{raw objective and evidence}
\rightarrow \text{capability-graph invention}
\rightarrow \text{mission compilation}\\
&\rightarrow \text{weakness-driven goals}
\rightarrow \text{missing-tool acquisition}
\rightarrow \text{verified consequences}\\
&\rightarrow \text{transfer, repair, and persistent capability update}.
\end{split}
\end{equation}
The result does not establish unrestricted self-education or AGI.  File-type
priors, safe adapter operations, semantic authorities, and the evaluated source
families remain engineered.  Arbitrary package installation, credentials,
purchasing, deployment, and unrestricted generated-code execution remain
separate authority classes and were not enabled.

The mission, curriculum, and bounded execution-acquisition track is therefore no
longer the principal architectural blocker.  The resumed critical path is:
\begin{enumerate}
    \item full-stack failure localization and private self-repair;
    \item multi-day heterogeneous continual learning over changing missions;
    \item stronger open natural-language comprehension and dialogue;
    \item broader delayed software, public-data, and physical consequences; and
    \item native consolidation of verified cross-domain trajectories while
    preserving HexCore as the final authority.
\end{enumerate}
